\section{Kerangka Berpikir}

Kerangka berpikir penelitian merangkum hubungan antara premis masalah, sintesis kajian pustaka, rancangan solusi, hipotesis, dan aktivitas metodologi. 
Alur tersebut disajikan pada Gambar \ref{fig:kerangka-berpikir}.

\begin{figure}[H]
    \centering
    \includegraphics[width=0.9\textwidth]{kerangka-berpikir.jpg}
    \caption{Kerangka Berpikir Penelitian}
    \label{fig:kerangka-berpikir}
\end{figure}

Gambar \ref{fig:kerangka-berpikir} menunjukkan bahwa Premis 1–4 menjadi dasar pembentukan tiga prinsip rancangan CoDA-EDA. 
Prinsip tersebut kemudian diterjemahkan ke dalam Hipotesis H1–H4 dan diuji secara empiris melalui enam aktivitas DSRM \parencite{peffers_design_2007}, dengan luaran berupa algoritma yang tervalidasi pada perangkat \textit{edge-IoT} nyata. 
Kerangka ini menjadi acuan dalam perancangan algoritma, instrumen evaluasi, dan desain eksperimen.